\documentclass{article}\usepackage[T1]{fontenc}\begin{document}
The existing NativeBlocks.Nested.flatten preserves and reflects cartesian arrows over the source, including lifts of each specified base map. Every arrow into the flattened image of a specified nested word has a unique incoming lift when its source word is allowed to vary. This is strictly stronger than faithfulness and does not assert fullness. The same theorem holds after attaching an arbitrary category-valued presheaf as terminal payload. Weighted flattening compares independently chosen reindexings by a natural isomorphism, with identity, composition, and mixed base/payload diagrams on strict source fibres.

The fixed-length binary lemma uses the earlier common base/grade object and hom universes b,w. The nested flattening is exactly the length-tagged span map of PR 207: the base and grades have object universe u and hom universe max(u,v), and the endpoint data are split. The terminal payload is any ordinary category-valued presheaf in independent object and hom universes; no closure condition for iterating that payload is used. The base coherence and mixed equations allow separate source and target cleavages. Normality and literal equality of independently selected lifts are not assumed.

Recover the distinction in lines 929--953 of projects/contextual-composition/foundations/relative-construction/delivery/One\_relative\_construction\_for\_context\_content\_and\_effects.tex at 458ff36a4798e432dee5947afc402f9903a47f14: native multiplication is a map of spans, while weighted comparison with the packed monad is further structure. Inspect the earlier NativeBlocks.blocks strict isomorphism, binary tag splitting and Nested.extend/collapse/flatten. Read Ragged207Review.lean: flatten is not full, so the previous SpanPayload.cartesian\_iff argument for fully faithful comparisons cannot apply. The surviving local control surface is a fixed codomain with a variable source, not a fixed source-target hom-set. Recheck this argument if length-changing arrows or identifications of block boundaries are introduced.

BlockCartesian.Incoming is the dependent sum of each domain with its actual morphism to a fixed object. UniqueIncoming means bijectivity on these sums. It is closed under composition, disjoint-sum descent and strict base change of the prefix in a matched pair. The latter retains the right arrow and its matching square and recovers the left source and arrow together. Each fixed pair of lengths is glued by an isomorphism of categories. Combining that fact with tagged descent and then inducing over the actual snoc recursion gives flattenIncoming. The factorization property for relative cartesian arrows is equivalent to unique factorization of a pair of incoming arrows; the bijections transport that property in both directions. This also proves a pointwise equivalence for the destination projection if an arrow there is cartesian: no fibration structure on that destination is claimed.

A cartesian cleavage on the nested source is built independently: lift the prefix, then lift the suffix over the destination image of the prefix lift and the matching transport. The universal factor in the strict-match category keeps the two component maps. No flattening or pullback of a target choice is used in this recursion. On a weighted total category incoming arrows decompose into an incoming grade arrow and the dependent payload source and arrow; changing the grade by flattening leaves that remaining dependent data in place. Hence All.weightedStrong preserves and reflects the strongly cartesian property over a specified base map. All.weightedFace supplies its comparison and cartesian factorization; All.unitLaw, compLaw and mixedLaw refer to this face. mixedLaw compares the two paths with the actual change-of-payload faces and requires no invertibility of the payload transformation. These are statements on strict source fibres, not a monad comparison.

MovingLift.lean uses natural-number contexts, additive extensions and a terminal presheaf indexed by finite positions. The independently constructed lift for a word of inner lengths 1,0,2 moves its source from 5 back to 2 and its endpoint from 20 to 17. Flatten sends the lift to a strongly cartesian but noninvertible arrow over 2 to 5. A vertical increase in a native grade fails the strongly cartesian condition even after flattening. At terminal position 16 the pulled payload has value 24; a nonidentity payload transformation and a nonidentity base arrow instantiate the mixed diagram. One incoming native arrow has distinct unique lifts for the empty outer word and the single empty block as specified codomains. Choice.lean gives integers as native grade arrows. In the length-one summand, choices have lifts 2 and 7 over the same identity; the comparison on that actual flattened fibre has label -5. The strict-preservation shortcut would predict 0 and is refuted.

Strict preservation of an independently chosen cleavage is not a consequence and is false for the displayed nonnormal choices. No claim about compatible normal choices for strict algebras is made. Three-level flattening associativity, action on maps of spans, a monad comparison identifying weighted flattening with PackedFree.join via evaluation, and strict contextad algebra recovery remain open. In particular, the base-composition coherence of the weighted flattening face is not associativity of three nested words. Issue \#105 stays open.
\end{document}
